% Nouveau EGA V (cote 25): archive cote 25, site page 4 (PDF page 5).
% typescript. Draft compared: none.
% AI-transcribed by Claude Opus 5.5 (claude-opus-5-5, adaptive thinking, effort high), Message Batches API, batch pilot (msgbatch_0189Nhf1xB615YWnztubtLfc), 2026-11-12; not yet checked by eye.
\scanpage{25/4}
\mspage{\hand{1}}
IV 20 \quad \uline{Sections hyperplanes et projections \var{\deleted{linéaires}}{}{deletion}{ega25_25_004-v0001}\var{\inserted{coniques}}{coniques}{insertion}{ega25_25_004-v0002}}.

20.1. \quad \uline{Préliminaires \var{\over{de}{et}}{et}{writing over}{ega25_25_004-v0003} notations}.

Soient S un \quad préschéma , $\underline{E}$ un Module loc libre de type fini sur S,
$\overset{\hand{\vee}}{\underline{E}}$ son dual. Nous noterons par P=$\underline{P}(\underline{E})$ le fibré projectif défini par $\underline{E}$,
par $\overset{\hand{\vee}}{P}$ celui défini par $\overset{\hand{\vee}}{\underline{E}}$ . Nous\unsure{n}appelerons \ $\overset{\hand{\vee}}{P}$ \ le \uline{schéma des hyperplans}
de P, cette terminologie se justifiant de la façon suivante. \ Soit
$\hand{\xi}$ une section de $\overset{\hand{\vee}}{P}$ sur S, qui est donc\hand{/}déterminée par un module quotient
inversible L de $\overset{\hand{\vee}}{\underline{E}}$ . On en conclut un module quotient inversible
$\underline{L}_P$ de$(\overset{\hand{\vee}}{\underline{E}})_P$ = $(\underline{E}_P)^{\hand{\vee}}$ , d'autre part on dispose d'un module quotient inversibl
$\underline{O}_P(1)$ de \ $\underline{E}_P$ . Passant \ aux duals, on peut envisager \ $\underline{L}_P^{-1}$ resp. $\underline{O}_P(-1)$
comme des sous-modules inversibles (loc facteurs directs) de \ $\underline{E}_P$ resp.
de $(\underline{E}_P)^{\hand{\vee}}$ , \ et l'accouplement \ $\underline{E}_P \otimes \overset{\hand{\vee}}{\underline{E}}_P \hand{\longrightarrow} \underline{O}_P$ \ définit \ donc un
accouplement naturel
(x) \[ \underline{O}_P(-1)\otimes\underline{L}_P^{-1} \hand{\to} \underline{O}_P \]
ou encore un homomorphisme \ transposé
\var{\margin{$\|$}}{\placedmargin{$\|$}}{marginal note placed}{ega25_25_004-v0004}
(xx) \[ \underline{O}_P \hand{\longrightarrow} \underline{O}_P(1)\otimes\underline{L}_P \ = \ \underline{D}_P(1) \ , \]
i.e. une section de $\underline{L}_P(1)$, canoniquement déterminée par $\hand{\xi}$ . Le''diviseur''
de cette section, \ i.e. le sous-schéma\var{\inserted{fermé}}{ fermé}{insertion}{ega25_25_004-v0005} $H_{\hand{\xi}}$ de \ P défini par l'Idéal
image de (x), \ est \ appelé l'\uline{hyperplan dans P} \ défini par \ l'élément
$\hand{\xi} \hand{\in} \overset{\hand{\vee}}{P}(S)$ . \ Il peut encore se décrire en notant que localement sur S \var{\xout{la}}{}{typewriter strike-out}{ega25_25_004-v0006}
\var{\xout{\unsure{donnée de $\check{\xi}$ équivaut}}}{}{typewriter strike-out}{ega25_25_004-v0007} \quad $\hand{\xi}$ \ est donné par une section \ $\phi$ de \ $\underline{E}$ , \var{\xout{de}}{}{typewriter strike-out}{ega25_25_004-v0008} telle
\var{\xout{\illegible{1}}}{}{typewriter strike-out}{ega25_25_004-v0009} que $\phi(s) \neq 0$ pour tout s \ ( $\phi$ étant déterminé par $\hand{\xi}$ \ modulo
multiplication par une section inversible de $\underline{O}_S$ ); comme $\underline{E}$ = $p_{\hand{*}}(\underline{O}_P(1))$
( p: P$\hand{\to}$S étant la projection), \ $\phi$ \ peut être considéré comme une section
de \ $\underline{O}_P(1)$, don\var{\over{c}{t}}{t}{writing over}{ega25_25_004-v0010} le diviseur n'est autre que $H_{\hand{\xi}}$. \ Bien entendu, si on
\var{\inserted[typed]{\uline{considére}}}{\uline{considére}}{insertion}{ega25_25_004-v0011}\var{\xout{\illegible{1}}}{}{typewriter strike-out}{ega25_25_004-v0012} $\underline{L}^{-1}$ \ comme un sous-\var{\xout{\illegible{1}}}{}{typewriter strike-out}{ega25_25_004-v0013} module inversible de \ $\underline{E}$ locale
ment facteur direct dans $\underline{E}$, alors la correspondance entre $\hand{\xi}$ \ (i.e. $\underline{L}$ \hand{\unsure{$\xi$,}}

\mspage{\hand{4}}
